\title{OPT-IML: Scaling Language Model Instruction Meta Learning through the Lens of Generalization}

\begin{abstract}
Recent work has shown that fine-tuning large pre-trained language models on a collection of tasks described via instructions, a.k.a. instruction-tuning, improves their zero and few-shot generalization to unseen tasks. However, there is a limited understanding of the performance trade-offs of different decisions made during the instruction-tuning process. These decisions include the scale and diversity of the instruction-tuning benchmark, different task sampling strategies, fine-tuning with and without demonstrations, training using specialized datasets for reasoning and dialogue, and finally, the fine-tuning objectives themselves. In this paper, we characterize the effect of instruction-tuning decisions on downstream task performance when scaling both model and benchmark sizes. To this end, we create OPT-IML Bench: a large benchmark for Instruction Meta-Learning (IML) of 2000 NLP tasks consolidated into task categories from 8 existing benchmarks, and prepare an evaluation framework to measure three types of model generalizations: to tasks from fully held-out categories, to held-out tasks from seen categories, and to held-out instances from seen tasks. Through the lens of this framework, we first present insights about instruction-tuning decisions as applied to OPT-30B and further exploit these insights to train OPT-IML 30B and 175B, which are instruction-tuned versions of OPT. OPT-IML demonstrates all three generalization abilities at both scales on four different evaluation benchmarks with diverse tasks and input formats -- PromptSource, FLAN, Super-NaturalInstructions, and UnifiedSKG. Not only does it significantly outperform OPT on all benchmarks but is also highly competitive with existing models fine-tuned on each specific benchmark. We release OPT-IML at both scales, together with the OPT-IML Bench evaluation framework.
\end{abstract}

\section{Introduction}

Instruction fine-tuning is shown \citep{wei2021flan,sanh2021t0,chung2022scaling} to significantly improve the zero- and few-shot performance of large pretrained LMs (LLM). It involves fine-tuning LLMs on collections of NLP tasks using instructional style input formats. Successful instruction-tuning of LLMs depends on a number of aspects such as the objectives used for fine-tuning, the distribution and diversity of the fine-tuning tasks, the inclusion of specialized datasets related to reasoning and dialogue, fine-tuning with demonstrations, and also, the comprehensiveness of the evaluation framework. In this paper, we develop an extensive large-scale fine-tuning and evaluation framework of 2000 NLP tasks (which we call OPT-IML Bench) and use it to characterize the tradeoffs of different decisions relating to instruction meta-learning (IML) on the OPT models \citep{zhang2022opt}. We exploit insights gathered from this process, to train OPT-IML 30B and 175B, instruction-tuned versions of OPT.

There are a growing number of large meta-datasets of NLP tasks such as Super-NaturalInstructions \citep{wang2022niv2}, FLAN \citep{wei2021flan} and PromptSource \citep{sanh2021t0}. Recent instruction-tuning work has demonstrated success using these individual benchmarks and their combinations~\citep{chung2022flanpalm}, with a general recommendation for scaling up the number of tasks. We follow this recommendation by consolidating 8 meta-datasets into a large collection of 1,991 NLP tasks containing instructions with multiple prompts and grouping them into more than 100 task categories such as Question Answering and Sentiment Analysis (Figure \ref{fig:flow}). Furthermore, we transform this collection into an evaluation framework for comprehensively evaluating large-scale instruction-tuned models across three levels of generalization: 1) model performance on tasks from fully held-out task categories not used for tuning, as in prior work \citep{wei2021flan,sanh2021t0}, and additionally, 2) performance on unseen tasks from categories seen during instruction-tuning, and, 3) performance on held-out instances of tasks seen during tuning. The former two settings evaluates the cross-task generalization of instruction-tuning while the last setting evaluates the generalization of supervised multi-task learning~\citep{mccann2018decathlon}.

We refer to the resulting instruction-tuning framework as OPT-IML Bench and illustrate its composition in Figure \ref{fig:flow} where the right hand side depicts evaluation categories, which can be completely disjoint, partially overlap, or completely overlap with the categories used for tuning on the left. Each category comprises datasets that can belong to multiple benchmarks and be associated with multiple prompts.

The effectiveness of instruction-tuning on LLMs depends on factors such as the diversity and distribution of tuning-tasks, the formatting of their prompts, and the objectives used for fine-tuning. Several recent works on instruction-tuning explore these factors by grouping tasks into categories and evaluating performance on tasks from completely held-out task categories \citep{sanh2021t0,wei2021flan,wang2022niv2}. Using our evaluation framework that considers multiple levels of generalization, we are able to comprehensively characterize the tradeoffs relating to these different factors when scaling up instruction-tuning to an aggregate of 8 different benchmarks. By instruction tuning OPT 30B \citep{zhang2022opt} on OPT-IML Bench, we outline the tradeoffs of dataset and benchmark sampling strategies during tuning, the scaling laws with respect to tasks and categories, the effects of approaches to incorporating task demonstrations into instruction-tuning based on \cite{min2021metaicl}, as well as instruction-tuning with specialized datasets that contain reasoning chains \citep{kojima2022large, wei2022cot} and dialogue \citep{shuster2022blenderbot}. These experiments can serve to establish best practices for large scale instruction-tuning of LLMs.

Given the insights gathered from our generalization experiments on OPT-IML bench, we train OPT-IML. OPT-IML significantly improves over its base pre-trained model at both 30B and 175B scales on four different instruction-tuning benchmarks: PromptSource \citep{sanh2021t0}, FLAN \citep{wei2021flan}, Super-NaturalInstructions \citep{wang2022niv2}, and UnifiedSKG \citep{xie2022unifiedskg}. Additionally, the OPT-IML models also perform competitively in comparison with each of the prior instruction-tuned models individually tuned on these benchmarks on both zero and few-shot performance. Recently, along similar lines as this work, \cite{chung2022flanpalm} achieve impressive gains on the challenging benchmarks of MMLU \citep{hendrycks2020measuring} and Big-Bench Hard \citep{suzgun2022challenging} by instruction-tuning PaLM \citep{chowdhery2022palm} and T5 \citep{raffel2020t5} on a scaled-up collection of 1.8K tasks. OPT-IML trained under similar settings still underperforms in comparison on these challenging benchmarks and we discuss this in Section \ref{sec:discussion}. Following OPT \citep{zhang2022opt}, we will responsibly share versions of OPT-IML at both scales, and also release our OPT-IML Bench evaluation framework to facilitate future work in this direction. 

\section{Scaling up Multi-task Benchmarks}
\label{sec:scaled-multi-task-benchmark}

To characterize the effects of extreme task scaling on instruction tuning, we build on recent task collections such as Super-NaturalInstructions \citep{wang2022niv2} and PromptSource \citep{sanh2021t0}, and aggregate 8 such collections to create the OPT-IML Benchmark for massive instruction fine-tuning and evaluation over diverse task categories, instruction types and prompting setups (Table \ref{tab:benchmark_compilation}). 

For the remainder of this paper, we use the terms task and dataset interchangeably; each task/dataset can be instantiated using multiple prompt templates. We refer to the original data from which the tasks are created as a data source; multiple tasks can be created from the same data source (e.g. question answering and question rewriting).  A benchmark comprises multiple tasks, where each task belongs to a single task category/cluster. 

\subsection{Task Curation}
\label{sec:unifying_benchmarks}
We expand the Super-NaturalInstructions benchmark of 1600+ tasks by~\citet{wang2022niv2} with the task collections from multiple existing work on \emph{instruction-tuning}: FLAN~\citep{wei2021flan}, \textsc{T0}~\citep{sanh2021t0}; \emph{prompt crowdsourcing}: PromptSource~\citep{bach2022promptsource};

\emph{cross-task transfer studies}: ExMix~\citep{aribandi2021ext5}, \textsc{T5}~\citep{raffel2020t5}, CrossFit~\citep{ye2021crossfit}; and \emph{area-specific task consolidation}: Structured Knowledge Grounding~\citep{xie2022unifiedskg}, Dialogue~\citep{shuster2022blenderbot} and Chain-of-thought Reasoning\footnote{We use 14 Chain-of-thought Reasoning datasets which form a superset of those used by~\cite{chung2022flanpalm} (Appendix~\ref{subsec:reasoning-benchmark}).}~\citep{chung2022flanpalm}. The curation process of all these benchmarks can be found in Appendix~\ref{sec:data_curation}. \\

There is a significant overlap between the datasets in these benchmarks. For example, popular datasets such as SQuAD v1/v2~\citep{rajpurkar2016squad,rajpurkar2018squadv2} appear in almost all benchmarks. In addition, while Super-NaturalInstructions, PromptSource, FLAN and Chain-of-thought Reasoning contain long-form human-written instructions or reasoning chains, the rest of the benchmarks are designed for multi-task learning and the prompt templates often only consist of short field or task prefixes (e.g. ``question:'', ``label:''). Therefore, we only kept tasks from the CrossFit, ExMix and \textsc{T5} collections that do not appear in any other benchmarks. Since we're exploring a large number of tasks, we take maximally 100k examples (at random) per task from all benchmarks except FLAN, where we take maximally 30k examples per task following the same practice as~\cite{wei2021flan}.

\subsection{Benchmark Consolidation}

\paragraph{Instruction schema.} Each benchmark adopts different instruction and language styles. In Table~\ref{tab:prompt_variation}, we broadly classify their instructions into two categories: dataset-level and instance-level. \emph{Dataset-level instructions} define the overall task and may include auxiliary information such as positive/negative examples and explanations. The model is expected to learn the definition of the task based on this and apply the knowledge to each example coming after it. \emph{Instance-level instructions} are templates to be instantiated for each example individually and is sometimes designed in the cloze-style to solicit the desired output for the example.

We cast all tasks across the benchmarks we collect into the bipartite prompt formulation that include ``instructions'' and ``output'' segments (Table~\ref{tab:prompt_variation}). 
For CrossFit, ExMix and \textsc{T5}, since the original benchmarks do not provide natural language instructions, we manually write a simple instruction sentence for each of the included tasks and use them at the instance level. For example, the instructions for the GPT-2 Deepfake Detection task~\citep{radfordgpt2output} in ExMix reads ``Is the following text produced by GPT-2?''.

\paragraph{Task categorization.} 

We categorize the tasks under the conventional NLP categories following the practice of previous work~\citep{wei2021flan,sanh2021t0,wang2022niv2,ye2021crossfit}. Such grouping offers a convenient scaffold to study the generalization of models cross- and within categories.

We primarily follow the 76-category taxonomy defined by Super-NaturalInstructions. The other benchmarks also provide their own task clusters. We perform a coarse unification of the task clusters manually, e.g. merging ``hate speech detection'' with ``toxic language detection''. Besides this, benchmarks such as CrossFit and PromptSource adopt a finer-grained task categorization compared to Super-NaturalInstructions, e.g. CrossFit identifies multiple sub-classes of Question Answering. In such cases, we adopt the more coarse-grained assignment of Super-NaturalInstructions. This results in a single-level taxonomy with over 100 task categories (Table~\ref{tab:benchmark_compilation}). 

\subsection{Creating Benchmark Splits}
\label{sec:evaluation_splits}
\paragraph{Train, validation and test splits.} We split the set of all tasks in a way that allows us to perform massive instruction fine-tuning and evaluate the resulting model with respect to three levels of generalization. First, we hold out several task categories to evaluate model generalization to \emph{new categories of tasks}. Second, we select a subset of the remaining categories as partially held-out categories.\footnote{We manually examined the full task collection to eliminate false negatives for the held-out and partially held-out categories.} We divide the datasets in these categories into train and evaluation and use them to test model generalization to \emph{new datasets from seen task categories}. We select the fully and partially held-out categories by largely staying consistent with previous instruction fine-tuning work~\citep{wang2022niv2,wei2021flan,sanh2021t0} to allow direct comparison. Finally, for a subset of the training tasks, we hold out the validation and test sets from the original data release, and use them to test model generalization in the standard multi-task learning setting, i.e. \emph{new examples from seen tasks}. We reserve 35 evaluation tasks spanning 9 task categories from the evaluation tasks as the validation set\footnote{We also added the validation split of the Measuring Massive Multitask Language Understanding benchmark~\cite{hendryckstest2021mmlu} in our experiments in \S\ref{sec:ablation}.}, and use them to characterize the tradeoffs of different instruction-tuning strategies in \S\ref{sec:ablation}. The details of our validation tasks including their evaluation metrics are shown in Table~\ref{tab:validation_tasks}.

\paragraph{Task de-duplication.} We make sure that the train and evaluation tasks do not overlap on the data source they were created from, to prevent leakage\footnote{This condition is maintained for our partially held-out evaluation tasks as well.}, following the practice of~\cite{wang2022niv2}. For each pair of train and eval tasks, we compute the fraction of examples that have any 13-gram overlap between the instantiated sequences from those examples. We manually examine every pair where more than $1\%$ of the eval set overlaps with the training set ($\sim$14,000 pairs) to confirm whether tuning on the train task can unfairly benefit the eval task, and decide either to remove the train or the eval task in confirmed cases. The task pairs that share a broad contextual resource such as Wikipedia but otherwise contain unrelated output labels are retained. Table~\ref{tab:benchmark_compilation} shows the statistics of our task splits.

\subsection{Task Prompt Construction}

Each example in the zero-shot setting is formatted using the bipartite instruction scheme as described in Section \ref{sec:unifying_benchmarks}. We insert a delimiter between the instructions and the output if the instructions do not end with a ``:''. Similar to~\cite{chung2022flanpalm}, for each example we randomly sample a delimiter from a small set\footnote{The set includes ``$\backslash$nAnswer:'', `` Answer:'', ``$\backslash$nA:'', `` A:'', ``$\backslash$nOutput:'', `` Output:'', ``$\backslash$nanswer:'', ``$\backslash$output:''.} to mitigate overfitting. 

For few-shot prompts, we place the demonstration examples between the task descriptions and the target example for benchmarks that adopt task-level instructions such as Super-NaturalInstructions, and before the task example for benchmarks that adopt instance-level instructions such as FLAN and PromptSource.

Examples of prompts for each of the tasks can be found in Appendix \ref{sec:templates-and-example-prompts}.

The FLAN and PromptSource benchmarks contain multiple manually-written templates per task. To further increase task diversity, some templates in these benchmarks altered the original task semantics (e.g. ``question answering'' $\rightarrow$ ``question generation''). We manually examined all task templates in these benchmarks and removed the templates that altered the original task semantics to refine our task categories.

\section{Instruction Fine-tuning}
\label{sec:modeling}

We use the OPT-IML Bench presented in Section \ref{sec:scaled-multi-task-benchmark} to fine-tune OPT \citep{zhang2022opt}, a suite of open-source decoder-only transformer language models released in scales from 125M to 175B parameters that performs similar to GPT-3 \citep{brown2020gpt} on a collection of standard NLP tasks. OPT is trained on 180B unique tokens from a combination of the datasets used in RoBERTa \citep{liu2019roberta}, the Pile \citep{gao2020pile}, and PushShift.io Reddit \citep {baumgartner2020pushshift,roller2020recipes} using a next-word prediction objective. 

We describe the process of instruction-tuning OPT at the scales of 30B and 175B in this section. 

\subsection{Fine-tuning Objective}

We finetune OPT in a manner similar to pre-training using a next-word prediction objective conditioned on all previous tokens as context. However, we separate the training sequence into a source context sequence and a target sequence and only include loss terms from the tokens in the target sequence (label-loss). We treat the task instructions and inputs as source tokens and the label tokens as target tokens. Formally, for a fine-tuning dataset $\mathcal{D}$ comprising source instances $s_i$ and their corresponding target tokens $t_i=\{t_{ij}\}$, a pre-trained model with parameters $\theta$ is fine-tuned to minimize the following loss over the target tokens conditioned on the source tokens and previously seen target tokens.

\begin{align}
    \mathcal{L}(\mathcal{D}; \theta) = - \sum_i \sum_j \log p_{\theta}(t_{ij}|s_i,t_{i,<j})
\end{align}

We minimize this loss across all datasets in our OPT-IML Bench by mixing examples from different datasets based on their sizes and proportions assigned to the benchmarks they come from (more details in Section \ref{sec:ablation}).

\subsection{Packing and Document Attention}
\label{subsec:packing}
In order to utilize the maximum sequence length for computational efficiency, we pack multiple examples (source and target) together as a sequence of 2048 tokens \citep{raffel2020t5}, separated by \texttt{<eos>} tokens. One consequence of packing is that the tokens belonging to one example can attend to tokens from previously packed examples in the same sequence. To mitigate this, we use document attention masking i.e. we modify the token attention mask in causal LMs to attend only to the tokens that are part of the same example, rather than all the previous tokens in the sequence. This changes the attention mask from a triangular to a block triangular mask and improves both stability and performance in our experiments. 

\subsection{Fine-tuning Hyperparameters}
\label{subsec:fthp}

We fine-tune all 30B models on 64 40GB A100s, and 175B models on 128 40GB A100s. Following OPT, we use Fully Sharded Data Parallel \citep{artetxe2021efficient} and the Megatron-LM Tensor Parallelism \citep{shoeybi2019megatron}. We inherit most model hyper-parameters for each model scale following OPT. We pack our training examples into sequences of length $2048$, left-truncating examples that overflow. We use Adam \citep{kingma2014adam} with 32-bit state with $(\beta_1, \beta_2) = (0.9, 0.95)$, linearly warming up the learning rate for $60$ steps to the maximum, followed by linearly decaying it to 0. We conduct preliminary experiments to select learning rates from $\{1e^{-5}, 3e^{-5}, 5e^{-5}, 6e^{-5}\}$ and per-GPU batch sizes from \{2, 4, 8\} using our validation split from \S\ref{sec:scaled-multi-task-benchmark}. The resulting hyperparameters are listed in Table \ref{tab:hyp}. We use a dropout of 0.1 (including embedding dropout) and clip gradient norms to 1.0, and use dynamic loss scaling to prevent underflows \citep{micikevicius2018mixed}. During fine-tuning, our models saw approximately 2 billion tokens, which is only 0.6\% of the pre-training budget of OPT (Table \ref{tab:hyp}).

\section{What Matters for Instruction Fine-tuning?}
\label{sec:ablation}

Recent works have explored a number of instruction fine-tuning techniques to optimize the performance of the resulting model on specific kinds of downstream tasks, and also to improve their robustness against variations in prompts, instruction styles and prompting setups. Using an OPT 30B model with the basic hyper-parameter settings chosen in \S\ref{subsec:fthp},

we run experiments to characterize the effects of dataset proportions, number of tasks and diversity, using pre-training, dialogue, and reasoning datasets, and training using demonstrations, on instruction-tuning with respect to our three levels of model generalization: \emph{fully held-out}, \emph{partially held-out} and \emph{fully supervised}. We aggregate performance along several dimensions such as clusters, and benchmarks to determine the best settings. 

\subsection{Experimental Setup}

The goal of our experimental setup is first, to characterize the effects of a multitude of factors related to the fine-tuning process, on instruction-tuning performance, and second, to use these findings to effectively instruction-tune OPT models. The factors that we experiment with are 1) the composition of the fine-tuning dataset mixture, 2) the number and diversity of the tasks used for fine-tuning, 3) using additional datasets relating to pre-training, reasoning and dialogue as part of the fine-tuning mix, and 4) different ways of fine-tuning with demonstrations. 

\paragraph{Prompt construction details.} To compile our train data, we merged all prompt data for a task with $N$ examples and randomly take $N$ prompts from the pool such that the training task distribution is kept the same regardless of how many prompts are given for the tasks. We merged the prompts for each task in a similar manner in our validation set, and randomly sample a maximum of 250 prompts per task to report the validation results. For our test tasks, we keep all prompt variations and all examples.

\paragraph{Generalization levels.} Starting with a baseline instruction-tuned model, we independently characterize the effect of each factor, by tuning models with several variations of that factor and evaluating the models on the tasks from our validation split from Section \ref{sec:scaled-multi-task-benchmark}, separated into three generalization levels: a) tasks from clusters not included in training (Fully Held-out), b) tasks unseen during training but from seen clusters (Partially Supervised), and c) tasks seen during training (Fully Supervised). An instruction-tuning setting is desirable if it improves performance on fully held-out and partially supervised tasks without sacrificing performance on fully supervised tasks. We use average performance across all three generalization levels on both 0-shot and 5-shot settings on the validation/test sets of the tasks in the validation split to determine the best settings for each factor. 

\paragraph{Decoding.} 
\label{sec:decoding}
Our evaluation data comprises tasks with answer candidates (of which one is correct), as well as tasks with multiple gold reference sequences. For the former set of tasks, we use rank classification similar to \cite{gpt3}, where we score each candidate based on their likelihood and output the highest-scoring candidate as the answer. This candidate is used to compute accuracy on the task. For tasks without candidates, we perform greedy decoding until an \texttt{<eos>} token is predicted or a maximum of \textit{N=256} tokens are generated. Based on the generated sequence and the references, we then compute either Exact-match or Rouge-L F1 scores.

\paragraph{Model selection.} For all experiments, we first aggregate results separately for 0-shot and 5-shot across task subtypes. For example, pro and anti versions of type 1 and type 2 Winobias \citep{zhao2018gender} tasks from PromptSource, and all 57 subtasks of MMLU \citep{hendrycks2020measuring}, would be aggregated to get per task performance. If the same task exists across multiple benchmarks, we then average performance across benchmarks as well. We then compute 0-shot and 5-shot averages of all tasks within a category (or benchmark depending on the experiment), and finally, compute a combined average of all 0 and 5-shot scores of each category (or benchmark), which we use for model selection. 

We tune each model for 4000 steps and evaluate on our validation split on both 0-shot and 5-shot settings, using 250 examples from each task for compute-efficiency. As described in Section \ref{sec:scaled-multi-task-benchmark}, our validation splits for each task include a mix of multiple prompts for FLAN and PromptSource. All but four validation tasks are generation-style tasks (where we report Rouge-L F1). We compute accuracy based on scoring for the remaining tasks and aggregate them together with Rouge-L for presentation purposes. We refer to Table~\ref{tab:validation_tasks} in the Appendix for full details about the tasks in our validation split. 

\subsection{Effects of varying task mixing-rate maximum}
\label{subsec:eps}

Prior work \citep{raffel2020t5,wei2021flan} typically uses example-proportional sampling and builds batches by sampling from datasets proportional to their sizes, while enforcing a maximum size parameter (EPS) to prevent large datasets from overwhelming the batch. To understand how this maximum mixing rate (EPS) affects performance across the different generalization levels, we perform experiments with EPS $\in \{128, 256, 512, 1024, 2048, 4096, 8192, 16384, 10^6\}$ and report results in Table \ref{tab:eps}. An EPS of 512 causes 97\% datasets to hit their maximum, while an EPS of 8192 causes 16\% datasets to hit their maximum. We also experiment without using EPS i.e. EPS=100K. 

Overall, we find that while EPS is important to instruction-tuning i.e. on average all models that use EPS outperform the model without it, after a certain threshold i.e. less than 4096 in our case, there is minimal variation in performance across all generalization levels. While based on the highest average performance, we choose 4096 (also corresponds to 50\% of the dataset lengths being capped) for our other experiments and the final OPT-IML models, we find that all values below 4096 also perform quite well, with EPS=128 closely matching 4096. Also note that changing EPS implicitly changes the proportion of fine-tuning data from each benchmark, which we control for explicitly in the next Section.

\begin{comment}

\end{comment}

\subsection{Effects of varying benchmark proportions}
\label{subsec:prop}

In Section \ref{sec:scaled-multi-task-benchmark}, we describe the multiple tasks and prompt repositories \citep{sanh2021t0,wang2022niv2,wei2021flan,ye2021crossfit,aribandi2021ext5} that we unify to massively scale the number of tasks used for instruction-tuning. However, using multiple benchmarks for training, together with only example-proportional sampling, results in benchmarks with more tasks overwhelming the batch composition. For example, in our benchmark, 71\% of training examples would come from SuperNatInst, with 18\% from PromptSource, and only 5\% from FLAN. Since each benchmark is associated with a specific task format, this can bias the resulting model towards certain input-output formats. We vary the proportions of different benchmarks to evaluate their effect on downstream task performance on our three generalization levels and present results in Table \ref{tab:prop}. For this experiment, we compare models based on their aggregate performance on each benchmark instead of task category, since we would like to choose the parameters that perform well on a maximum number of benchmarks.

First, we look at performance improvements within the same benchmark where the proportions were changed. As we increase the proportion of FLAN from 5\% to 25\%, its performance improves significantly on both the fully-held out and the partially held-out generalization levels, with no notable improvement on the fully-supervised tasks. SuperNatInst shows a similar trend on partially-supervised tasks, but surprisingly, not so much on fully held-out tasks. It is possible that the very specific input-output format of SuperNatInst makes it such that changing proportions of unrelated clusters provides no benefit to its fully held-out clusters. PromptSource is relatively unchanged on fully supervised clusters and partially supervised clusters, possibly owing to reaching performance saturation with even an 18\% proportion. However, it benefits with more proportion on the fully-held out clusters. 

Secondly, we also observe benchmarks complementing each other. For example, the highest accuracy on fully held-out FLAN i.e. 88.8/83.6\%, is achieved, not with the highest proportion of FLAN, but with improving the proportions of PromptSource and Crossfit. Similarly, the highest generation performance on fully-held out PromptSource of 79.7/83.5\% is achieved with 25\% PS, and not with 45\% PS proportions. We also observe certain tradeoffs, for example, the best proportions for FLAN and PromptSource result in a sharp drop in performance on reasoning datasets, and vice versa. Finally, setting Crossfit, Exmix, T5 and Unified-SKG proportions to 0 results in the worst model, demonstrating the benefits of using a diverse set of benchmarks for instruction-tuning. Based on average performance across benchmarks, ``2/1/27/40/27/1/2", ``7/1/35/25/28/2/2" and ``4/2/20/25/45/2/2" performed the best and we choose the last one as the proportion for our final OPT-IML models. Despite our choice, instruction-tuned models with different end-goals (for example, producing reasoning chains) would benefit from choosing differently. We also explore methods to improve performance on reasoning datasets in Section \ref{sec:cot}.

\subsection{Effects of Scaling Tasks or Categories}
\label{subsec:scaling-tasks-clusters}

Previous work has shown that scaling the number of training tasks or clusters improves the overall performance of the model on the fully held-out generalization setting~\citep{wei2021flan,wang2022niv2}. We study effects along similar axes but with more generalization settings such as fully held-out, partially supervised, and fully supervised tasks/categories. We use cluster/category interchangeably in this section. For the task scaling study, we randomly sample 16, 64, 256, and 1024 sets of tasks such that smaller sets are subset of bigger sets, and fully supervised tasks are always selected. Figure~\ref{fig:tasks_abalation} (full results in Appendix Table~\ref{tab:scaling-tasks}) presents these task scaling studies on the three generalization levels, aggregated at the cluster-level for both 0 and 5-shot performance. 

We observe that both fully held-out and partially supervised tasks get the most improvements with the increase in the number of training tasks. Interestingly, fully supervised tasks' performance remains unchanged even when more relevant tasks are seen from the fully supervised tasks' clusters, as we increase the training tasks. In the fully held-out setting, \emph{Cause Effect Classification} and \emph{Word Analogy} clusters see the biggest improvements in zero-shot and few-shot, respectively. On the partially supervised, \emph{Question Answering} and \emph{Toxic Language Detection} clusters see the biggest improvements on both zero-shot and few-shot. 

For the cluster scaling study, we order the clusters based on the decreasing order of the number of tasks present in each cluster and select the first 4, 16, 64, and 93 (i.e., all) clusters. Additionally, we make sure that Question Answering, Summarization, and Dialogue Generation clusters are always represented since our fully supervised validation tasks belong to these three clusters. Figure~\ref{fig:num_clusters} (full results in Appendix Table~\ref{tab:scaling-clusters}) presents the corresponding results on all three generalization levels for both zero-shot and few-shot settings. We observe that as we increase the training clusters, the performance on fully supervised tasks either stay the same or slightly drop in the few-shot setting. On the fully held-out and partially supervised levels, the results on the zero-shot settings improve an increase in the number of clusters and the results are a bit mixed for the few-shot setting, but overall they tend to decrease with cluster scaling. Note that the first 4 clusters already cover 673 tasks (clusters belonging to the fully supervised setting have a lot of tasks). Hence, the model starts with strong performance, which might lead to the mixed results that we observe. Based on these experiments, we use all tasks and clusters to train our final OPT-IML models.

\subsection{Effects of Pre-training during Instruction-Tuning}
\label{subsec:pretrain-data}

We observe that using pre-training style updates on entire sequences during fine-tuning can make training more stable, so we explore the performance effects of using pre-training data on our three generalization levels. Table \ref{tab:ablation-data-examples} shows an example used in the pre-training style updates. Following \cite{shuster2022blenderbot}, we use the last shard of the corpus used to train OPT~\citep{zhang2022opt} as our pre-training data for fine-tuning, since it is seen only once during the pre-training stage of OPT. We experiment with adding pre-training data by proportion in the increasing amounts of 1\%, 5\%, 10\%, and 50\%, and present results for the 5-shot setting, aggregated by task category, in Figure~\ref{fig:pretrain-data} (full 0 and 5-shot results in Appendix Table~\ref{tab:pretrain-data}).  

Overall, for the fully held-out and partially supervised generalization levels, we observe that the model improves while adding pre-training data for up to 10\% and then starts deteriorating after that. We also observe that using more pre-training data leads to better Rouge-L F1 scores but lower accuracy scores, partly owing to the influence of pre-training data on the remaining proportions of generation vs. classification tasks. Based on the average scores across generalization levels (see Appendix Table~\ref{tab:pretrain-data}), we choose to include 5\% pre-training data in instruction-tuning our OPT-IML models.

\subsection{Effects of Adding Reasoning Datasets}
\label{sec:cot}
Recent work \citep{wei2022cot,kojima2022large} has illustrated improvements in the performance of LLMs on reasoning tasks, when prompted to generate a reasoning chain in natural language before generating the answer. Based on these findings, we attempt to explicitly fine-tune LLMs to perform reasoning by compiling a set of 14 reasoning datasets (see Appendix \ref{subsec:reasoning-benchmark} for a list of these datasets), where the output includes a rationale before the answer and by including these datasets during instruction-tuning. This set includes the 9 datasets used by \cite{chung2022flanpalm} in their CoT category as well as some additional datasets. Each dataset has a single prompt that uses an instruction, that explicitly asks the model to generate a reasoning chain \citep{kojima2022large}, followed by examples in the few-shot setting that illustrate how the reasoning chain should be produced before the answer. We show an example with such a prompt in Table \ref{tab:ablation-data-examples}. Using benchmark proportions of ``2/1/27/40/27/1/2" as a baseline (see Section \ref{subsec:prop}), we experiment with adding 1\%, 2\%, and 4\% proportions of reasoning data (by reducing the proportion of the highest proportion benchmark i.e. SuperNatInst), and present results for the 5-shot setting in Figure~\ref{fig:cot} (full 0 and 5-shot results in Appendix Table \ref{tab:cot}) by generalization level and task category. 

We see a substantial performance improvement on the 2/14 held-out validation reasoning tasks (Rouge-L from 12.2\% to 31.6\%) when we instruction-tune with reasoning datasets, but alongside, we also see improvements on other held-out task categories such as Cause-Effect, Stereotype Detection, Toxicity Detection, and Word Analogy. Furthermore, adding 1\% reasoning data results in the largest gains overall, beyond which, the gains start to reduce on MMLU, Cause-Effect Accuracy, Toxicity, and Dialogue (averaged over 0 and 5-shot). On the other hand, the Summarization cluster (only 0-shot, see Appendix) continues to benefit from higher proportions of reasoning data. Based on average performance across categories and generalization levels, we use 1\% reasoning data for our final OPT-IML models.

\begin{comment}

\end{comment}

\subsection{Effects of Adding Dialogue Datasets}
\label{sec:ablation_dialogue}
We experiment with adding dialogues as auxiliary fine-tuning data to test if it can improve the LM's ability to respond to directional input and understand referential expressions. Another goal is to evaluate if this approach can induce chat-bot behaviors~\citep{shuster2022blenderbot} and make the resulting models more conversational. Using a subset of dialogue datasets\footnote{Appendix~\ref{subset:dialogue-datasets} list the dialogue datasets we used in this experiments.} used for training BlenderBot 3~\citep{shuster2022blenderbot}, we process the dialogues into sequences of turns separated by a single newline token (see Table \ref{tab:ablation-data-examples} for an example). The data consists of  320,543 unique dialogues and we fine-tune the model to predict the entire dialogue sequence. We set the proportion of the included dialogue data to be 0.5\% and present 0 and 5-shot results by task category and generalization level on our validation split in  Table~\ref{tab:add_dialogue}.

We observe that adding even just 0.5\% of the aforementioned dialogue data lowers 0-shot performance while 5-shot performance remains unchanged. Specifically, 0-shot performance suffers mainly on stereotype detection and word analogy. On examining model predictions on these categories, we found that they are primarily generation tasks whose references are either a single word or a short piece of text with a specific format (for example, a pair of phrases from the original input that refer to each other). Training with BB3 data weakened the model's ability to conform to the required format.\footnote{On one hand this behavior demonstrates a weakened instruction-following ability for the underlying model. On the other hand, it exposes a caveat in measuring model performance on tasks with instructions -- model performance on a specific task category is often the result of multiple factors and underperforming on a particular task category may not offer a useful atomic diagnosis. As in our case, we found the model to perform worse on stereotype detection tasks because it cannot parse the required output format, not because it is a more biased model.} It also significantly lowered the 5-shot performance of toxicity detection. An error analysis revealed a similar problem i.e. the model tends to perform worse on tasks that require generating a special set of decision words rather than simply generating ``yes'' or ``no''. Owing to severe model degeneration on these tasks, we do not add dialogue data while tuning OPT-IML.

\subsection{Effects of Meta-Training for In-Context Learning}
\label{subsec:meta-icl}

Recent work has shown that fine-tuning language models with demonstration examples in the instructions improves their ability to learn from the examples in context~\citep{min2021metaicl,wang2022niv2,chung2022flanpalm}. Both~\cite{min2021metaicl} and~\cite{wang2022niv2} experimented with the setup where a constant number of $k$ demonstration examples are added to each training example. The models are evaluated with the same number of $k$ demonstration examples during inference. ~\cite{chung2022flanpalm} used a mixture of data with and without exemplars. However, the proportion of each type of data used and how many exemplars were included are not clear. 

We attempt to train models that are better in-context few-shot learners, and also robust to the number of demonstration examples used during inference time.\footnote{In preliminary experiments, we found models trained with $k$ exemplars tend to perform worse when a different number of exemplars is used during inference time. Table~\ref{tab:niv2} shows a similar effect on the T$k$-\textsc{Instruct} models.} We experiment with a simple way of creating training examples that include varying numbers of demonstration examples. For each example $e$, we sample $k$ from a distribution $\mathcal{D}$ with cap\footnote{For any $k>K$, we set $k=K$.} $K$, and  randomly select $k$ other examples $E_d=\{e_1,\hdots,e_k\}, e_i \neq e$, from the train set, if $k > 0$. We add $E_d$ as the demonstration examples in $e$'s prompt, where the examples are separated by a special token \texttt{[SEP]}. For benchmarks with task-level instructions such as Super-NaturalInstructions, we place the demonstration examples before $e$ and after the instruction field; for benchmarks with instance-level instructions such as FLAN and PromptSource, we place the demonstration examples before $e$. 

Because the demonstration examples significantly increase the prompt lengths, including too many few-shot training examples often leads to worse performance and reduced learning stability, owing to sparsity in the loss and lower batch diversity. As a result, we choose $\mathcal{D}$ to be the Zipf distribution\footnote{\url{https://docs.scipy.org/doc/scipy/reference/generated/scipy.stats.zipf.html}}, which can be heavily tilted towards $k=0$. We train MetaICL models with different $\mathcal{D}$'s by adjusting the shape parameter $a$ of the Zipf distribution. When $a=4$, $92.5\%$ of the examples are zero-shot examples; and when $a=2$, $67.1\%$ of the examples are zero-shot examples. We set $K=5$ and use three consecutive newline tokens as \texttt{[SEP]} following~\cite{min2021metaicl}. 

\paragraph{MetaICL with suffix loss.} To further address the loss sparsity problem, we also experiment with a variation of the original MetaICL loss, illustrated in Figure~\ref{fig:metaicl_loss}. Given an example with instructions and exemplars, rather than training the model to produce the target label, we train the model to produce the target label of the first exemplar followed by the complete sequences of the remaining exemplars. This effectively turns the demonstration examples into training examples as well, and mitigates the loss sparsity problem given it is now spread over more tokens. 

\paragraph{Performance degradation on generation tasks.} We present validation set results for instruction-tuning with different settings for MetaICL, aggregated by generalization level and task category under both 0 and 5-shot settings, in Table~\ref{tab:metaicl_ablation}. We observe that adding MetaICL training leads to worse performance in both 0-shot and 5-shot setups in most cases, while MetaICL with the suffix loss outperforms regular MetaICL, especially in the 0-shot setup. Further examination of per-category performance reveals that while MetaICL models show reasonable improvements in multiple 5-shot evaluations, the 5-shot performances on Stereotype Detection and Word Analogy degrade significantly. An error analysis reveals a similar problem as in  \S\ref{sec:ablation_dialogue} -- the MetaICL models tend to lose the ability to strictly follow the output pattern in the presence of in-context exemplars. In addition, the standard MetaICL loss significantly hurts reasoning tasks. The resulting models tend to generate short answers despite the presence of reasoning chains in the in-context learning examples. Further investigation reveals that the model could be over-fitting to the demonstration separators and modifying them at inference time can significantly mitigate these problems (Table~\ref{tab:metaicl_double_newline_eval}).\footnote{Fine-tuning with random demonstration separators may effectively mitigate these issues and we will investigate this approach.} Interestingly, MetaICL degrades performance only for generation tasks, but is overall beneficial for scoring based classification tasks such as MMLU. However, owing to severe output degeneration in the regular setting, we decide to not use MetaICL to train our OPT-IML models.

\section{OPT-IML Models} 
\label{sec:results}

Using the best settings for instruction tuning from our experiments in Section \ref{sec:ablation}, we instruction tune OPT 30B and 175B to create OPT-IML 30B and 175B models. Specifically, we choose the best values for EPS and benchmark proportions, include all tasks in the training split, add 1\% datasets with reasoning chains, and 5\% data from the OPT pre-training corpus, and choose to leave out training with demonstrations i.e. MetaICL, as well as dialogue datasets. We tune OPT-IML 30B for 4000 steps, while we tune OPT-IML 175B for double the number of steps with half the batch size (for purposes of memory efficiency). Based on periodic validation set metrics, we decide to use the last checkpoint as the final model.

We evaluate our OPT-IML models on the OPT evaluation tasks as well as on four multi-task benchmarks from prior work \citep{wei2021flan,sanh2021t0,wang2022niv2,xie2022unifiedskg,zhang2022opt} in both zero and 5-shot settings, directly comparing them to individual benchmark specific instruction-tuned models released by prior work. Thus, we compare with baseline OPT models on the evaluation sets used by OPT, with FLAN-137B on the evaluation sets of FLAN \citep{wei2021flan}, with T0pp 11B on the evaluation sets from PromptSource \citep{sanh2021t0}, with Tk-Instruct 11B on the evaluation sets from Super-NaturalInstructions \citep{wang2022niv2}, and on joint modeling of text with code/structs on three tasks from the UnifiedSKG \citep{xie2022unifiedskg} benchmark. We examine these results in the following sections, and find that OPT-IML outperforms OPT on all benchmarks and is competitive with the individual benchmark specific instruction-tuned models on both zero- and few-shot performance.

\subsection{OPT Evaluations}
\label{sec:opt}

We evaluate OPT-IML on a subset of 14 standard NLP tasks reported by OPT \citep{zhang2022opt} on zero and few shot settings at 30B and 175B scales, using the same prompts released by OPT (a single prompt per task). All these tasks are classification-style tasks with multiple candidates, so similar to OPT, we use the candidate with the highest likelihood as the model prediction and report accuracies in Table \ref{tab:opt}. Additionally, all these tasks are held-out during training, some from our fully-held out categories and some from our partially held-out categories. For the few-shot setting, we use the same examples and number of shots used by OPT i.e. 32-shots, but truncated to fit within the model's maximum sequence length.  

On average, OPT-IML improves over OPT with approximately 6-7\% on 0-shot accuracy at both 30B and 175B model scales. For 32-shot accuracy, we see significant improvements on the 30B model, and milder improvements on 175B. While the improvements are significant for certain tasks such as RTE, WSC, BoolQ, ARC, CB, and WiC, our instruction-tuning does not improve performance for other tasks such as StoryCloze, PIQA, Winograd,  and Winogrande. Some of these latter results are specific to the prompts used by OPT. For example, we observe improvements on StoryCloze and Winogrande, when evaluated on a collection of prompt templates as part of PromptSource in Section \ref{subsec:eval-promptsource}. One reason for this is that OPT prompts were originally adopted from GPT-3 \citep{brown2020gpt} and have gone through a process of prompt engineering for optimal performance, while FLAN and PromptSource evaluate accuracies as averages using a diverse collection of prompts, including sub-optimal prompts. Thus, an advantage of instruction-tuning for these tasks can be to improve model robustness and reduce the need for prompt engineering.

\subsection{Evaluations on PromptSource}
\label{subsec:eval-promptsource}
\cite{sanh2021t0} fine-tune an LM adapted version of T5 11B\citep{raffel2020t5,lester2021power} on ~50 datasets from PromptSource (called T0) and evaluate on a set of 11 held-out tasks which are part of their 4 fully held-out categories. Each task is associated with multiple prompt templates, contributed by the research community with the help of their prompting tool. Since all these tasks are also part of held-out categories in OPT-IML, we use a similar evaluation setup, with some additional tasks as well. Most tasks are classification tasks where we score candidates based on likelihood and report accuracy, with the exception of Blended Skill Talk, which is a generation task where we report Rouge-L F1 scores. Since each task uses multiple prompts, we report metrics averaged across prompts under 0-shot and 5-shot settings in Table \ref{tab:promptsource}. 

Some of the prompts gathered in PromptSource are for an inverted version of the task. For example, the inverted task for QA is question generation. We do not train or evaluate using these prompts, since they are problematic when tasks are assigned to categories. We compare OPT-IML with the T0-original-task-only model which corresponds to our held-out setup (\cite{sanh2021t0} also release T0p and T0pp trained with additional tasks), and is also trained only on prompts that adhere to the original task. 

OPT-IML 175B matches the zero-shot performance of T0-original-task (11b) and outperforms it significantly on 5-shot performance. While both models were not trained on demonstrations, causal LMs like OPT demonstrate stronger generalization to the few-shot setting than encoder-decoder models like T0, and the latter could benefit from MetaICL training to improve its few-shot performance, as explored by \cite{chung2022flanpalm}. Similarly, on the Blended Skill Talk generation task, T0 underperforms causal LMs, which could be attributed to the large scale of the tuning data for OPT-IML, or may highlight a difficulty for encoder-decoder models to generalize to new generation tasks. At both scales, OPT-IML outperforms baseline OPT models on almost every task except Crows Pairs. As described in Section \ref{sec:opt}, this evaluation uses multiple prompts per task and rewards models that are more robust to the input prompts. Additionally, note that OPT-IML 30B outperforms baseline OPT 175B on average, demonstrating that instruction-tuning can be a way to make smaller-scale resource-efficient models more competitive. 

Following \cite{sanh2021t0}, we also evaluate on the WinoGender Schemas \citep{rudinger2018winogender} cast as a textual entailment task \citep{poliak2018dnc}, which measures the extent of gender bias in LLMs, and find that instruction-tuning vastly improves accuracy on this task. Finally, we evaluate on Crows Pairs \citep{nangia2020crows} formulated as a boolean QA task about whether a sentence illustrates a stereotype or not (using a single prompt), and see a deterioration in performance on OPT-IML 175B over OPT, but not on the 30B model. It is possible that other formulations of this task, for example, predicting which sentence is a stereotype, may show different trends. Note that these two tasks are not from held-out clusters, so there may be other training datasets that are beneficial. 

\subsection{Evaluations on FLAN}

Together with the FLAN instruction-tuning benchmark comprising 62 datasets, which we include in OPT-IML Bench, \cite{wei2021flan} also use it to instruction-tune Lamda-PT \citep{thoppilan2022lamda}, a 137B causal LM trained on 1.5T words of public dialog data and web text. They evaluate instruction-tuning using FLAN-137B on fully held-out task categories, by using a leave-one-out strategy i.e. they tune on all other categories, thus producing a different model to evaluate each test category. This presents an opportunity for evaluating OPT-IML models on the same task categories to assess the improvements that can be achieved by scaling up the instruction tuning benchmark to 1500 tasks using a single instruction-tuned model. 

We evaluate our OPT-IML models on a subset of tasks used by FLAN-137B, and based on our splits, some tasks are from fully-held out categories (ANLI, CB, MNLI, RTE, SNLI, WNLI, Winogrande, WSC), while the remaining are from partially held-out categories (BoolQ, OpenBookQA, ARC). All these tasks use a classification style with answer candidates, which we evaluate by scoring based on likelihood, and we report zero-shot and few-shot accuracies in Table \ref{tab:flan}. Note that each task is associated with 7-10 templates, and we report average accuracy across all templates. Some templates invert the task (for example, QA becomes question generation), and we do not evaluate on these templates. Also, while FLAN-137B uses a different number of shots for each task for their few-shot evaluation, we report 5-shot results for all tasks. 

We find that instruction-tuning significantly improves performance over baseline OPT models at 30B as well as 175B scales on each of the 15 tasks individually. While \cite{wei2021flan} found instruction-tuning to hurt fully-held out tasks at 8B and lower scales, but showing emergent behavior at a scale of 66B parameters and beyond, our experiments do not show this emergent behavior i.e. both 30B and 175B OPT-IML models achieve more than 20\% average improvement over the respective untuned models under 0-shot and few-shot settings. Additionally, our 30B OPT-IML model outperforms the 175B base OPT model by 20\% on 0-shot and 12\% on 5-shot, illustrating that instruction-tuned models at lower scales can be strong resource-efficient alternatives to larger untuned models. Compared with FLAN-137B, OPT-IML 175B performs competitively on 5-shot performance, and yields an improvement of 3\% on average on 0-shot performance. Nevertheless, the various differences in experimental setup relating to the held-out clusters, model sizes and the number of pre-training tokens, make it difficult to definitively attribute these improvements to scaling up the instruction-tuning benchmark.

\subsection{Evaluations on Super-NaturalInstructions}

Different from the evaluations seen so far, Super-NaturalInstructions uses a strict instructional format (Section \ref{sec:scaled-multi-task-benchmark}), where a formal instruction block is provided at the start of the prompt, detailing option candidates and resolving task ambiguities, followed by multiple demonstrations, and can help assess the ability of our models to generalize to different instruction formats. \cite{wang2022niv2} subdivide the SuperNatInst benchmark into training and held-out categories, and train Tk-Instruct 3B and 11B, which are instruction-tuned versions of LM-adapted T5 models. They evaluate Tk-Instruct on 12 categories representing 154 tasks for fully held-out generalization. Of these 12 categories, Textual Entailment, Coreference Resolution and Dialogue Act Recognition are fully held-out in our evaluation framework. We evaluate OPT-IML on these three categories in 0-shot, 2-shot and 5-shot settings and report Rouge-L F1 scores in Table \ref{tab:niv2}. These three categories comprise 44 tasks and we evaluate on the top-100 examples from these tasks following \cite{wang2022niv2}, with each task using a single prompt. In all cases, we generate a maximum of 256 tokens for each test example. For comparison, we also re-evaluate Tk-Instruct 11B on these clusters under the same evaluation framework. We use the version of Tk-Instruct 11B that performs best overall i.e. the version trained with instructions + 2 positive demonstrations and no negative demonstrations. 

Since Tk-Instruct is trained and evaluated under a 2-shot setting, we additionally report results on the 2-shot setting for this evaluation. First, OPT-IML models outperform baseline OPT models on each cluster at both scales, under 0-shot and all few-shot settings and once again we observe that an instruction tuned 30B model outperforms an untuned 175B model. Also, while both OPT 30B and 175B perform comparably at all shots, the instruction-tuned version of 175B vastly outperforms OPT-IML 30B, showing that larger models can benefit more from instruction tuning. Note that different from the Textual Entailment and other tasks from previous evaluations, all tasks here are evaluated under the generation setting (as opposed to scoring), which makes it significantly harder for untuned models. OPT-IML 175B outperforms Tk-Instruct 11B on 0-shot formats despite the former being tuned on a mixed-set of diverse formats from multiple benchmarks, whilst the latter being specifically tuned for this benchmark. The trend is reversed for the 2-shot and 5-shot settings where Tk-Instruct outperforms OPT-IML. Here, OPT-IML shows uniform performance under both settings whereas Tk-Instruct is heavily biased towards the 2-shot setting for which it was trained. Thus, the performance of Tk-Instruct drops from 65.3 to 58.4, from 2-shot to 5-shot. 

\subsection{Evaluations on UnifiedSKG}

UnifiedSKG \citep{xie2022unifiedskg} is a collection of 21 tasks related to Structured Knowledge Grounding with heterogeneous inputs such as databases, dialogue states, SQL queries, etc., which we include in OPT-IML Bench purposefully to equip the model with capabilities for handling structured knowledge. To evaluate these capabilities, we compare OPT-IML models with baseline OPT on three UnifiedSKG tasks formatted as text-to-text: DART \citep{nan2020dart}, which is a held-out data-to-text task for transforming data triples to text, Spider \citep{yu2018spider}, a SQL query generation task given a database and an input query, and fully supervised in our framework, and MultiWoZ \citep{budzianowski2018multiwoz}, is a held-out dialogue state tracking task. All three tasks are generation tasks where we decode 256 tokens before stopping and report Rouge-L F1 scores under 0-shot and 5-shot settings in Table \ref{tab:uskg}. 

On Spider, which is a fully supervised setting, OPT-IML models retain high performance close to a Rouge-L F1 score of 85  despite the presence of numerous other tasks in the instruction-tuning mix. On DART, OPT-IML shows modest gains in the 5-shot setting, but significantly outperforms OPT models on the zero-shot setting, with OPT-IML 30B outperforming OPT 175B. MultiWoZ, on the other hand shows significant deterioration with instruction tuning at both model scales.

\section{Discussion and Limitations}
\label{sec:discussion}
In the previous section, we demonstrated on multiple evaluation benchmarks that effectively instruction-tuned models can obtain significant improvements over untuned models on both zero- and few-shot settings. We achieved this by first scaling up the instruction-tuning datasets to encompass 8 large collections of NLP tasks, which we transform into an evaluation framework that tests \emph{three} levels of model generalization on downstream tasks. 

Using this framework, we characterized the tradeoffs of different factors on instruction tuning such as 1) the number and diversity of input tasks, 2) the distribution of different tasks and instruction styles, 3) the inclusion of specialized datasets relating to reasoning chains and dialogue, and 4) fine-tuning with demonstrations. This exploration helped us choose the best settings to instruction tune OPT-IML models at 30B and 175B scales, which perform competitively on an extensive set of benchmarks. 

In this section, we report additional results on instruction fine-tuning using our full task collection and discuss the limitations of our current approach.

\subsection{Evaluations on MMLU, BBH and RAFT}

While we transform our massively scaled instruction-tuning benchmark into an evaluation framework to study instruction-tuning techniques, recently \cite{chung2022flanpalm} also scaled up instruction fine-tuning up to 1,836 tasks from 4 benchmarks using the PaLM models~\citep{chowdhery2022palm} up to 540B and T5 models~\citep{raffel2020t5} up to 11B\footnote{Our work started concurrently.}.

The resulting models, namely the FLAN-PaLM and FLAN-T5 series were evaluated on several challenging language model benchmarks including MMLU~\citep{hendryckstest2021mmlu}, and Big Bench Hard (BBH)~\citep{srivastava2022beyond}. In order to establish the performance of OPT-IML in a similar setting (and additionally, on RAFT \citep{alex2021raft}), we instruction-tune OPT 30B and 175B on our entire benchmark of 1,991 tasks, which we call OPT-IML-Max.

We use option scoring for the two classification benchmarks MMLU and RAFT, and generation with Exact Match for BBH. We evaluate on the test sets for MMLU and BBH and on the evaluation split for RAFT released by the HELM benchmark~\citep{liang2022holistic}. We report these results in Table \ref{tab:bbh-mmlu} together with other large pre-trained and instruction-tuned models. Additionally, we also train and present results for OPT-IML-Max at the 1.3B scale (using the same settings as OPT-IML-Max 30B). On all three datasets,  OPT-IML-Max outperforms its untuned counterparts at all scales (except 1.3B on BBH). While, OPT-IML-Max is competitive with FLAN-T5 11B on RAFT, its performance lags behind FLAN-T5, FLAN-PaLM and the family of instruction-tuned GPT-3 models (*-davinci-*) on MMLU and BBH. 

 While the scale of the instruction-tuning benchmark remains similar across these models, there are many other underlying differences. There is a large variation with respect to the number of tokens used to train the respective underlying pre-training models. For example, T5 is trained on 1T tokens, FLAN-PaLM on 800B and OPT on 180B. There are also differences relating to the composition of the pre-training data and the respective modeling architectures. \cite{chowdhery2022palm} find that encoder-decoder models can fine-tune more effectively than decoder only models at similar scales, and massively scaling up decoder-only models can make them more competitive. Finally, there are also differences in the fine-tuning algorithms used, for example, some of the OpenAI davinci models use RLHF \citep{christiano2017deep} on feedback signals gathered from their API in addition to supervised fine-tuning. While we found that using Meta-ICL (\S\ref{sec:ablation}) did not yield a holistically better model and did not include it in our final models, they yielded 2-3\% improvements on MMLU and BBH. All these factors make it difficult to explain the gap in performance on these benchmarks, but nevertheless, these evaluations serve to establish the effects of our instruction tuning decisions with respect to OPT models on these challenging external benchmarks.

 \subsection{Limitations} 
 We use our evaluation framework to characterize the tradeoffs of various instruction-tuning variables on OPT 30B independently of each other. Although resource intensive to test, it is possible for these variables to interact with each other resulting in a different choice of the best tuning settings (for example, adding reasoning datasets may affect the choice of benchmark proportions). Furthermore, all tradeoffs studied on 30B instruction tuning may not show the same trends at larger scales. While we study instruction tuning tradeoffs using a comprehensive set of splits of fully held-out, partially supervised and fully supervised categories, choosing a different set of categories may result in prioritizing different decisions than those we took in this paper. Although we assign tasks to categories based on the underlying formats, such an assignment can be subjective and a different category assignment might change the optimal factors for instruction-tuning. For example, tasks that require different skills such as detecting toxicity can also be cast as textual entailment tasks. 

 \subsection{Responsible AI} While OPT-IML models outperform baseline OPT on an extensive set of evaluations (Section \ref{sec:results}), nevertheless, they are susceptible to the various risks associated with using large language models relating to factual correctness \citep{thoppilan2022lamda,brown2020gpt,chowdhery2022palm}, generation of toxic language \citep{gehman2020realtoxicityprompts} and enforcing stereotypes. While we release our OPT-IML models to proliferate future work on instruction-tuning and to improve the availability of large instruction-tuned causal LMs over 100B parameters, the use of these models should be accompanied with responsible best practices. 

\section{Related Work} Our work on fine-tuning large language models to follow instructions span multiple areas such as multi-task learning, prompting, and meta-training of in-context learning. We discuss these areas below within the scope that most closely relate to our work.

\paragraph{Instruction Tuning.} Language models are trained to predict the next token in a sequence with self-supervised learning~\citep{brown2020gpt,zhang2022opt,chowdhery2022palm}. Prompt engineering and in-context learning has become a dominant approach to leverage these models to solve many NLP tasks. In order to align these models to follow natural instructions and avoid prompt engineering, recent works have proposed instruction fine-tuning~\citep{ouyang2022gpt3instruct,wei2021flan,chung2022flanpalm,wang2022niv2}. Some of these works focus on fine-tuning the model on a wide range of tasks using human annotated prompts and feedback~\citep{ouyang2022gpt3instruct}, whereas the others focusing on supervised fine-tuning using academic benchmarks and datasets augmented with manually or automatically generated instructions~\citep{wang2022niv2,wei2021flan,sanh2021t0,Zhong2021AdaptingLM}. In our work, we focus on the second approach and consolidate a massive collection of publicly available datasets with instructions to finetune OPT. Concurrent to our work,~\citet{chung2022scaling} also proposes a similar instruction benchmark scaling approach to 1836 tasks from 4 benchmarks. While they focus on fine-tuning using the entire benchmark in order to push the limits of performance on several challenging held-out tasks that test the model's world knowledge and reasoning capabilities such as MMLU \citep{hendrycks2020measuring} and Big-Bench Hard (BBH) \citep{suzgun2022challenging}, we focus on characterizing the tradeoffs of various instruction-tuning decisions that can affect downstream performance. 

\paragraph{Prompting and Meta-Training} Zero- and few-shot learning (a.k.a. in-context learning) that leverages very few examples to solve any NLP task by effectively prompting the language models, is becoming a dominant paradigm in recent years~\citep{brown2020gpt}. Prompting involves modifying the input and output space of a given task that can effectively leverage the knowledge of the language model to solve it. Various approaches have proposed better prompting ways to improve generalization performance~\citep{wei2022cot,lu2021fantastically}. Furthermore, recent developments have shown ways to improve in-context learning (ICL) by meta-tuning language models to better adapt for ICL~\citep{min2022rethinking,min2021metaicl}. In our work, we leverage both the variants of prompts available from different benchmarks, as well as meta-training with demonstrations from a large pool of tasks, to study the effective settings for instruction-based fine-tuning that induce robustness against different prompting language and setups.

\paragraph{Learning to Reason.}

Despite the progress of in-context learning, state-of-the-art LLMs still struggle with reasoning tasks such as commonsense reasoning~\citep{west-etal-2022-symbolic}, and math word problems~\citep{hendrycksmath2021} which require arithmetic reasoning, etc. To solve these challenging tasks, recent work used different prompting methods which include a rationale with the final answer in the form of a scratchpad for arithmetic and logical reasoning~\citep{nye2021show}, provided chain-of-thought prompts in demonstrations~\citep{wei2022cot}, or added trigger phrases such as \textit{let's think step-by-step} to prompt models to generate explanations~\citep{kojima2022large}. In addition to changing prompts, \citet{chung2022scaling} integrated step-by-step explanations into the instruction tuning stage. Following \cite{chung2022scaling}, we further expand the set of reasoning datasets to 14 datasets and study the effects of different proportions of reasoning data on different held-out task clusters.

\paragraph{Multi-task Learning.} Instruction-based fine-tuning can be considered as a formulation of multi-task Learning (MTL). MTL is a popular paradigm that improves the generalization performances of a task when combined with related tasks by sharing common parameters or representations~\citep{caruana1997multitask,kumar2012learning}. MTL has been applied to many NLP scenarios in recent years primarily focusing on improving the performance on the training tasks or to new domains by leveraging the signal from related tasks~\citep{collobert2008unified,mccann2018decathlon,raffel2020t5,vu2020exploring}. In contrast, instruction-based fine-tuning allows us to improve the generalization performance to new tasks that are never seen during training. This is achieved by unifying all the tasks into a common format~\citep{kumar2016askmeanything,khashabi2020unifiedqa} via \emph{instructions}, and training them together by sharing all the weights of the model across all tasks. 

\paragraph{Continuous Learning.} Existing work also address continuous adaptation of language models by revisiting the instructions~\citep{yin2022contintin} or examples~\citep{scialom2022ct0} of previously learned tasks when fine-tuning with a new task to prevent catastrophic forgetting. The results show that LMs can be adapted effectively to new tasks without losing sight of the previously learned tasks. Other work enable the LM to perform new tasks via arithmetic combination of learned task vectors~\citep{ilharco2022taskvector} or soft prompts~\citep{anonymous2023progressiveprompts} patched to the base LM without changing its parameters. We focus on the (massively) multi-task adaptation setting by fine-tuning the LM with 2000 tasks at once. Continuously adapting the resulting model to new data, new tasks and new domain would be an interesting and important future direction.

\section{Conclusions}

Instruction-tuning of LLMs has emerged as an effective means to improve their zero and few-shot generalization abilities. We make three main contributions to instruction-tuning in this paper. First, we curate a large scale benchmark for instruction-tuning comprising 2000 NLP tasks from 8 dataset collections, annotated into task categories. We strategically produce evaluation splits on this benchmark to evaluate three different types of model generalization abilities: 1) fully-supervised performance, 2) performance on unseen tasks from seen task categories, and 3) performance on tasks from completely held-out categories. Second, using our evaluation suite, we establish  tradeoffs and best practices of many aspects of instruction-tuning, such as different sampling methods of fine-tuning tasks and categories, fine-tuning with task demonstrations, and fine-tuning with specialized datasets for reasoning and dialogue. Finally, using the best settings from our experiments, we train and release OPT-IML 30B and 175B instruction-tuned models based on OPT, that strongly outperform OPT on five evaluation benchmarks and are competitive with recent instruction-tuned models that are tuned on individual benchmarks.

\subsubsection*{Acknowledgments}
We would like to thank Stephen Roller, Susan Zhang, and Naman Goyal for help with fine-tuning OPT using the \texttt{metaseq} codebase and with our model release; Lili Yu for help with infrastructure and evaluations; Sewon Min for discussions related to meta-training for in-context learning; and Omer Levy, Timo Schick, and Scott Yih for helpful discussions related to instruction-tuning.

\end{document}
